\chapter{Mise en œuvre : extraits de code}\label{chap:code}

Ce chapitre présente des extraits du code source, choisis pour montrer comment chaque couche de l'architecture est réalisée : entités et relations Hibernate, configuration, accès aux données, service transactionnel, contrôleur, filtre et vue. Les extraits sont repris tels quels du dépôt ; les passages omis sont signalés par un commentaire.

\section{Entités et relations}

Les cinq entités portent des annotations JPA, à partir desquelles Hibernate génère les tables. L'entité \code{Commande} montre les trois types de correspondance utilisés : une énumération enregistrée sous forme de texte, une association \code{@ManyToOne} vers le client et une association \code{@OneToMany} vers les lignes. Cette dernière est une composition : avec \code{cascade = CascadeType.ALL}, enregistrer la commande enregistre aussi ses lignes.

\begin{lstlisting}[language=Java,caption={Entité \code{Commande} : colonnes, énumération et associations}]
@Entity
@Table(name = "commande")
public class Commande {
    // ...
    @Id
    @GeneratedValue(strategy = GenerationType.IDENTITY)
    private Long id;

    @Column(name = "date_commande", nullable = false)
    private LocalDateTime dateCommande = LocalDateTime.now();

    @Enumerated(EnumType.STRING)
    @Column(nullable = false, length = 20)
    private StatutCommande statut = StatutCommande.EN_ATTENTE;

    @Column(name = "adresse_livraison", nullable = false, length = 255)
    private String adresseLivraison;

    // RG11 : calculé par le service à partir des lignes
    @Column(name = "montant_total", nullable = false, precision = 10, scale = 2)
    private BigDecimal montantTotal = BigDecimal.ZERO;

    @ManyToOne(optional = false)
    @JoinColumn(name = "utilisateur_id", nullable = false)
    private Utilisateur utilisateur;

    // Composition : les lignes sont enregistrées et supprimées avec la commande
    @OneToMany(mappedBy = "commande", cascade = CascadeType.ALL, orphanRemoval = true)
    private List<LigneCommande> lignes = new ArrayList<>();
    // ... constructeurs
    // Maintient les deux côtés de l'association
    public void ajouterLigne(LigneCommande ligne) {
        ligne.setCommande(this);
        lignes.add(ligne);
    }
    // ... accesseurs
}
\end{lstlisting}

\begin{lstlisting}[language=Java,caption={Entité \code{LigneCommande} : classe d'association entre \code{Commande} et \code{Produit}}]
@Entity
@Table(name = "ligne_commande")
public class LigneCommande {

    @Id
    @GeneratedValue(strategy = GenerationType.IDENTITY)
    private Long id;

    @Column(nullable = false)
    private int quantite;

    // RG10 : copie du prix du produit au moment de la commande
    @Column(name = "prix_unitaire", nullable = false, precision = 10, scale = 2)
    private BigDecimal prixUnitaire;

    @ManyToOne(optional = false)
    @JoinColumn(name = "commande_id", nullable = false)
    private Commande commande;

    @ManyToOne(optional = false)
    @JoinColumn(name = "produit_id", nullable = false)
    private Produit produit;
    // ... constructeurs et accesseurs
}
\end{lstlisting}

\code{LigneCommande} porte la clé étrangère des deux associations (\code{commande\_id}, \code{produit\_id}) : c'est le côté propriétaire, d'où le \code{mappedBy = "commande"} du côté \code{Commande}. La méthode \code{ajouterLigne} met à jour les deux côtés de l'association en mémoire.

\section{Configuration d'Hibernate}

Le fichier \code{hibernate.cfg.xml} décrit la connexion à MySQL, demande la mise à jour du schéma au démarrage (\code{hbm2ddl.auto = update}) et déclare les classes annotées. La propriété \code{current\_session\_context\_class = thread} associe la session au thread courant, le temps d'une transaction : le service et les DAO partagent ainsi la même session et la même transaction.

\begin{lstlisting}[language={},caption={\code{hibernate.cfg.xml}}]
<hibernate-configuration>
    <session-factory>
        <!-- Connexion MySQL (identifiants de développement local) -->
        <property name="hibernate.connection.driver_class">com.mysql.cj.jdbc.Driver</property>
        <property name="hibernate.connection.url">jdbc:mysql://127.0.0.1:3307/boutique?characterEncoding=UTF-8</property>
        <property name="hibernate.connection.username">boutique</property>
        <property name="hibernate.connection.password">boutique</property>
        <property name="hibernate.connection.pool_size">10</property>

        <!-- Tables générées à partir des entités -->
        <property name="hibernate.hbm2ddl.auto">update</property>

        <!-- Une session par thread : le service ouvre la transaction, les DAO l'utilisent -->
        <property name="hibernate.current_session_context_class">thread</property>

        <property name="hibernate.show_sql">true</property>

        <mapping class="ma.ensa.boutique.model.Utilisateur"/>
        <mapping class="ma.ensa.boutique.model.Categorie"/>
        <mapping class="ma.ensa.boutique.model.Produit"/>
        <mapping class="ma.ensa.boutique.model.Commande"/>
        <mapping class="ma.ensa.boutique.model.LigneCommande"/>
    </session-factory>
</hibernate-configuration>
\end{lstlisting}

\section{Accès aux données : DAO générique}

Les opérations communes à toutes les entités sont écrites une seule fois dans \code{GenericDao<T>}. Chaque DAO (\code{ProduitDao}, \code{CommandeDao}, etc.) en hérite et n'ajoute que ses requêtes HQL. Les DAO n'ouvrent ni ne valident aucune transaction : ils utilisent la session courante.

\begin{lstlisting}[language=Java,caption={\code{GenericDao<T>} : opérations CRUD communes}]
// CRUD commun à toutes les entités ; la transaction est ouverte par le service
public abstract class GenericDao<T> {

    protected final Class<T> classe;

    protected GenericDao(Class<T> classe) {
        this.classe = classe;
    }

    protected Session session() {
        return HibernateUtil.getSessionFactory().getCurrentSession();
    }

    public void sauvegarder(T entite) {
        session().persist(entite);
    }

    public T modifier(T entite) {
        return session().merge(entite);
    }

    public void supprimer(T entite) {
        session().remove(entite);
    }

    // Identifiant absent ou mal formé : null, traité comme « introuvable » par le service
    public T trouverParId(Long id) {
        return id == null ? null : session().get(classe, id);
    }

    public List<T> listerTous() {
        return session().createQuery("from " + classe.getSimpleName(), classe).getResultList();
    }
}
\end{lstlisting}

\section{Service : passage de commande transactionnel}

Le service \code{CommandeService} applique les règles RG09 à RG12. La méthode \code{enTransaction} encadre le traitement : validation si tout se passe bien, annulation complète dès qu'une \code{RegleMetierException} est levée.

\begin{lstlisting}[language=Java,caption={\code{CommandeService.passerCommande}}]
public Commande passerCommande(Long utilisateurId, Panier panier, String adresseLivraison) {
    // RG09
    if (panier == null || panier.estVide()) {
        throw new RegleMetierException("Votre panier est vide.");
    }
    // ... contrôle de l'adresse de livraison
    // Les produits sont verrouillés par id croissant : deux commandes simultanées
    // prennent les verrous dans le même ordre et ne peuvent pas se bloquer mutuellement
    List<LignePanier> lignesPanier = new ArrayList<>(panier.getLignes());
    lignesPanier.sort(Comparator.comparing(ligne -> ligne.getProduit().getId()));

    return enTransaction(() -> {
        // L'utilisateur en session est détaché : on le recharge
        Utilisateur utilisateur = utilisateurDao.trouverParId(utilisateurId);
        if (utilisateur == null) {
            throw new RegleMetierException("Utilisateur introuvable.");
        }
        // Compte désactivé par un administrateur depuis la connexion (BF22)
        if (!utilisateur.isActif()) {
            throw new RegleMetierException("Votre compte est désactivé : vous ne pouvez plus commander.");
        }
        Commande commande = new Commande(utilisateur, adresse, StatutCommande.EN_ATTENTE);

        for (LignePanier lignePanier : lignesPanier) {
            // Le produit du panier est une copie qui peut être périmée : on relit prix, stock et état en base
            Produit produit = produitDao.trouverEtVerrouiller(lignePanier.getProduit().getId());
            // RG15 : produit supprimé ou désactivé depuis l'ajout au panier
            if (produit == null || !produit.isActif()) {
                throw new RegleMetierException("Le produit « " + lignePanier.getProduit().getNom()
                        + " » n'est plus disponible : retirez-le de votre panier.");
            }
            int quantite = lignePanier.getQuantite();
            // RG12 : une exception ici annule la commande et tous les stocks déjà décrémentés
            if (produit.getStock() < quantite) {
                throw new RegleMetierException("Stock insuffisant pour « " + produit.getNom() + " » ("
                        + produit.getStock() + " en stock, " + quantite + " demandé(s)).");
            }
            produit.setStock(produit.getStock() - quantite);
            // RG10 : prix actuel du produit copié dans la ligne
            commande.ajouterLigne(new LigneCommande(produit, quantite, produit.getPrix()));
        }

        // RG11
        BigDecimal total = BigDecimal.ZERO;
        for (LigneCommande ligne : commande.getLignes()) {
            total = total.add(ligne.getPrixUnitaire().multiply(BigDecimal.valueOf(ligne.getQuantite())));
        }
        commande.setMontantTotal(total);

        commandeDao.sauvegarder(commande);
        return commande;
    });
}
\end{lstlisting}

\begin{lstlisting}[language=Java,caption={\code{CommandeService.enTransaction} : validation ou annulation}]
private <R> R enTransaction(Supplier<R> traitement) {
    Session session = HibernateUtil.getSessionFactory().getCurrentSession();
    session.beginTransaction();
    try {
        R resultat = traitement.get();
        session.getTransaction().commit();
        return resultat;
    } catch (RuntimeException e) {
        session.getTransaction().rollback();
        throw e;
    }
}
\end{lstlisting}

\begin{lstlisting}[language=Java,caption={\code{ProduitDao.trouverEtVerrouiller} : lecture du produit avec verrou}]
// Passage de commande (RG12) : recharge le produit avec un verrou en écriture (select ... for update).
// Deux commandes simultanées sur le même produit sont ainsi traitées l'une après l'autre,
// ce qui empêche de vendre plus que le stock. Le verrou est libéré au commit ou au rollback.
// Requête HQL plutôt que get() : seule la ligne du produit est verrouillée, pas celle de sa catégorie.
public Produit trouverEtVerrouiller(Long id) {
    if (id == null) {
        return null;
    }
    return session()
            .createQuery("from Produit p where p.id = :id", Produit.class)
            .setParameter("id", id)
            .setHibernateLockMode(LockMode.PESSIMISTIC_WRITE)
            .uniqueResult();
}
\end{lstlisting}

\section{Contrôleur : servlet de commande}

La servlet lit les paramètres, appelle le service et choisit la suite : redirection après un succès (le message est déposé en session et affiché une seule fois), ou réaffichage du formulaire avec le message d'erreur. Elle ne contient aucune règle de gestion.

\begin{lstlisting}[language=Java,caption={\code{CommandeServlet.doPost}}]
@Override
protected void doPost(HttpServletRequest requete, HttpServletResponse reponse)
        throws ServletException, IOException {
    HttpSession session = requete.getSession();
    Utilisateur utilisateur = (Utilisateur) session.getAttribute("utilisateur");
    if ("annuler".equals(requete.getParameter("action"))) {
        annuler(requete, reponse, utilisateur);
        return;
    }
    Panier panier = (Panier) session.getAttribute("panier");
    String adresseLivraison = requete.getParameter("adresseLivraison");

    try {
        Commande commande = commandeService.passerCommande(utilisateur.getId(), panier, adresseLivraison);
        panier.vider();
        session.setAttribute("flashSucces", "Votre commande n° " + commande.getId()
                + " a bien été enregistrée. Montant total : "
                + String.format(Locale.FRANCE, "%,.2f", commande.getMontantTotal()) + " DH.");
        reponse.sendRedirect(requete.getContextPath() + "/client/commandes");
    } catch (RegleMetierException e) {
        // Le panier n'est pas modifié : le client corrige et confirme à nouveau
        requete.setAttribute("erreur", e.getMessage());
        afficherValidation(requete, reponse, adresseLivraison);
    }
}
\end{lstlisting}

\section{Filtre de contrôle d'accès}

\code{AuthFilter} s'applique à toutes les adresses \code{/client/*}. Il laisse passer la requête seulement si la session contient un utilisateur de rôle CLIENT ; sinon il redirige, avant même que la servlet ne soit appelée.

\begin{lstlisting}[language=Java,caption={\code{AuthFilter.doFilter}}]
// BNF02, RG08 : l'espace client est réservé aux utilisateurs connectés avec le rôle CLIENT (déclaré dans web.xml, après EncodingFilter)
public class AuthFilter implements Filter {

    @Override
    public void doFilter(ServletRequest req, ServletResponse rep, FilterChain chaine)
            throws IOException, ServletException {
        HttpServletRequest requete = (HttpServletRequest) req;
        HttpServletResponse reponse = (HttpServletResponse) rep;

        HttpSession session = requete.getSession(false);
        Utilisateur utilisateur = session == null ? null : (Utilisateur) session.getAttribute("utilisateur");

        if (utilisateur == null) {
            reponse.sendRedirect(requete.getContextPath() + "/connexion");
        } else if (utilisateur.getRole() == Role.ADMIN) {
            reponse.sendRedirect(requete.getContextPath() + "/admin/tableau-de-bord");
        } else {
            chaine.doFilter(requete, reponse);
        }
    }
}
\end{lstlisting}

\section{Vue : JSP avec JSTL et EL}

Les pages n'utilisent que les balises JSTL et le langage d'expression. L'extrait du panier parcourt les lignes avec \code{c:forEach}, affiche les textes avec \code{c:out} (qui échappe le HTML) et met en forme les prix avec \code{fmt:formatNumber}.

\begin{lstlisting}[language={},caption={\code{client/panier.jsp} : lignes du panier et total}]
    <c:forEach items="${panier.lignes}" var="ligne">
        <tr>
            <td>
                <a href="${ctx}/produit?id=${ligne.produit.id}"><c:out value="${ligne.produit.nom}"/></a>
            </td>
            <td class="text-end text-nowrap">
                <fmt:formatNumber value="${ligne.produit.prix}" minFractionDigits="2" maxFractionDigits="2"/> DH
            </td>
    <%-- ... formulaires de modification et de retrait --%>
    </c:forEach>
</tbody>
<tfoot>
    <tr>
        <th colspan="3" class="text-end">Total</th>
        <th class="text-end text-nowrap">
            <fmt:formatNumber value="${panier.total}" minFractionDigits="2" maxFractionDigits="2"/> DH
        </th>
        <th></th>
    </tr>
</tfoot>
\end{lstlisting}
